\documentclass[12pt]{article}
\usepackage[utf8]{inputenc}
\usepackage[T1]{fontenc}
\usepackage{hyperref}
\title{Negeri Sembilan Malay Language Resources for K-12 Curriculum Development}
\author{Compiled by Claude Sonnet 5.5 (claude-sonnet-5-5) for Andr\'as Kornai}
\date{}
\begin{document}
\maketitle

\section{Introduction}
Negeri Sembilan Malay (ISO 639-3 zmi; Glottocode nege1240) is spoken mainly in Negeri Sembilan, Malaysia, with pockets in Selangor, Malacca and Johor. Its speakers descend from Minangkabau settlers from Sumatra. Sources disagree on whether it is a Minangkabau variant or closer to Standard Malay. No standard dictionary, grammar or corpus specific to this variety was found in the searches.

\section{Primary Educational Resources}
No dedicated school materials were found. Standard Malay (Bahasa Malaysia) is the language of schooling in Malaysia.

\section{Linguistic and Reference Materials}
\begin{itemize}
\item Wikipedia, ``Negeri Sembilan Malay''. \url{https://en.wikipedia.org/wiki/Negeri_Sembilan_Malay}
\item Rufus Hendon, \emph{Phonology and Morphology of Ulu Muar Malay} (Yale, 1966), a study of a Negeri Sembilan dialect. The page found in the search is at \url{https://penangbookshelf.com/pd-phonology.cfm}
\item Springer volume \emph{Minangkabau and Negri Sembilan}, on the historical relationship between the two regions. \url{https://link.springer.com/book/9789400982000}
\end{itemize}

\section{Religious and Community Resources}
\begin{itemize}
\item Wikipedia, ``Minangkabau Malaysians''. \url{https://en.wikipedia.org/wiki/Minangkabau_Malaysians}
\end{itemize}

\section{Digital and Audio Resources}
\begin{itemize}
\item Global Recordings Network language page for zmi. It listed no recordings at the time of the search. \url{https://globalrecordings.net/en/language/zmi}
\end{itemize}

\section{Curriculum Development Context}
The material is mostly scholarly and scattered. A curriculum would have to draw on Malay-medium schooling and add local-dialect and cultural content, ideally documented with speaker communities. Minangkabau resources may be relevant but should not be assumed to apply directly.

\end{document}